% ================================================================
%  chapters/parts/ch05-a.tex
%  第 5 章  化学吸附和程序升温技术（分片 a：5.1—5.3）
%  源：S1（上册）第 5 章，书内 p.301—319 / PDF p.320—338
%  分片 b 覆盖 5.4—5.7（应用）
% ================================================================

\chapterstart{化学吸附和程序升温技术}%
{本章回答：怎样用化学吸附量得到活性位与金属分散度，三种吸附等温式各自假定了什么样的
表面，以及如何把一条程序升温脱附（TPD）谱线翻译成吸附量、脱附活化能和活性位数目。
导读按「原理—模型—动态方法」推进：5.1 给出化学吸附的热力学与速率基础，5.2 说明
Langmuir／Freundlich／Temkin 的适用边界，5.3 建立 TPD、TPR、TPO 的谱图判读规则。}
\label{ch:05}

\subsection{化学吸附的基本原理}

\begin{probout}
化学吸附测量的是\textbf{气体分子在固体表面以化学键形式发生的单层、选择性富集}，它给出的
不是几何信息，而是\textbf{活性位的种类、强度与数量}。本节建立三件事：吸附为什么自发、
吸附与脱附各以多快速率进行、以及这两者与催化活性的经验关系。
\end{probout}

\subsubsection{化学吸附过程的热力学讨论}

吸附是自发过程，$\Delta G_{\text{吸}}<0$；被吸附物失去全部或部分平动自由度，熵也降低，
故 $\Delta S_{\text{吸}}<0$。由 $\Delta G_{\text{吸}}=\Delta H_{\text{吸}}-T\Delta
S_{\text{吸}}$ 必有 $\Delta H_{\text{吸}}<0$：\textbf{吸附一律放热}，每个吸附过程有其特有
的吸附热 $q=-\Delta H_{\text{吸}}>0$\textsuperscript{\src{S1 §5.1.1，书内 p.304 /
PDF p.323}}。吸附可写为 $\ce{A + S <=> AS}$（$\ce{S}$ 为表面吸附位），其平衡常数与覆盖度
的关系即 5.2 节三种等温式的共同出发点。
此处已对照原页面图像（书内 p.304 / PDF p.323）：原书印「每一个吸附过程都有其特有的吸附热
$q$［$q=-\Delta H(\text{吸})$］」，吸附热即小写斜体 $q$，全页未出现 $Q$ 或 $q_{\text{ad}}$；
「吸」在原书用正体全角括号随行排（$\Delta G(\text{吸})$、$\Delta H(\text{吸})$、
$\Delta S(\text{吸})$），本稿改排为下标，属排版差异。原书同页印式 (5-1)
$\Delta G(\text{吸})=\Delta H(\text{吸})-T\Delta S(\text{吸})$、式 (5-2)
$\ce{A + S <=> AS}$（「式中，$\ce{S}$ 为表面吸附位；$\ce{A}$ 为吸附物」）、式 (5-3)
$K=[\ce{AS}]/([\ce{S}][\ce{A}])$，与本稿一致。

化学吸附的判据性特征是\textbf{由化学键引起、有电子转移或共享、吸附热大
（约 \qtyrange{50}{960}{kJ.mol^{-1}}，物理吸附仅 \qtyrange{10}{30}{kJ.mol^{-1}}）、
有选择性、单层、需同时加热与抽真空才能脱除}；物理吸附则由范德华力引起、无选择性、可多层、
抽真空即可除去。原书特别提醒\textbf{两者界线在许多情况下并不明显}，判据必须合并使用，
不能只看吸附量大小。

\subsubsection{化学吸附与催化活性的两条经验规则}

多相催化循环包括扩散、化学吸附、表面反应、脱附和反向扩散五个步骤，化学吸附是其中唯一
同时决定「活化什么」和「活化多少」的步骤。长期研究归纳出两条经验规则
\textsuperscript{\src{S1 §5.1.1，书内 p.303—304 / PDF p.322—323}}：

\paragraph{规则一（必要条件）}至少有一种反应物在催化剂表面上发生化学吸附，是产生催化活性
的必要条件；若对任何反应物都无化学吸附能力，就不可能催化该反应。

\paragraph{规则二（适度原则）}表面对反应物的吸附必须\textbf{适当}：多相催化需要的是较弱并且
快速的化学吸附，单位表面反应速率在相同覆盖度下与吸附强度成反比。

\begin{pitfall}
规则二常被误读成「吸附越强越好」。吸附太强会使表面被稳定中间体覆盖而无法脱附，活性反而
下降——即 Volcano 曲线的下降支。测到很大的吸附热只说明该位点可能难以周转，并不意味着
好催化剂。
\end{pitfall}

\subsubsection{吸附与脱附速率}

吸附速率由碰撞频率、吸附概率与空位分数共同决定，$r_{\text{ad}} = k_a\,C\,(1-\theta)$；
脱附速率为 $r_{\text{de}} = k_d\,\theta^{\,n}$（$\theta$ 为覆盖度，$n$ 为脱附级数），
速率常数按 Arrhenius 形式随温度变化：
\begin{equation}
k_d = A\exp\!\left(-\frac{E_d}{RT}\right)
\label{eq:ch05-kd}
\end{equation}
$E_d$ 为脱附活化能，是覆盖度 $\theta$ 的函数，$A$ 为脱附频率因子。平衡时令两速率相等即得
吸附等温式的一般形式。注意 $E_d$ 依赖 $\theta$ 这一点：\textbf{表面不均匀或存在横向相互
作用时，脱附活化能会随覆盖度漂移}，这正是 Temkin 型行为的物理来源。

\subsection{三种模型的吸附等温式}

\chapref{ch:01} 已从动力学与多层吸附假设分别推导了 Langmuir 与 BET 等温式，本节不再重复
推导，只回答一个操作性问题：\textbf{每一种等温式隐含了什么样的表面假设，以及在什么条件下
必须放弃 Langmuir 而改用 Freundlich 或 Temkin。}

\begin{principle}
\textbf{Langmuir}：$\theta = \dfrac{Kp}{1+Kp}$。假设表面\textbf{均匀}（各位点吸附热相同）、
\textbf{单层}吸附、被吸附分子之间\textbf{无相互作用}、吸附与脱附处于\textbf{动态平衡}；
推论是 $p/V$–$p$ 为直线，且饱和吸附量 $V_m$ 可达平台。

\textbf{Freundlich}：$V = K_F\,p^{1/n}$（$n>1$）。等价于假设表面具有\textbf{指数分布的
非均匀吸附热}；$\ln V$–$\ln p$ 为直线。它是纯经验式，\textbf{不给出饱和吸附量}，因此不能
用于测定金属分散度。

\textbf{Temkin}：$\theta = \dfrac{RT}{\alpha}\ln(K_0 p)$。假设吸附热随覆盖度\textbf{线性下降}
（$\Delta H_{\text{吸}} = \Delta H_0 - \alpha\theta$），即表面不均匀但下降方式规则；适用区间
为中等覆盖度（约 $0.2<\theta<0.8$），$\theta$–$\ln p$ 为直线。
\end{principle}

\begin{dataflow}
选模型的判据不是「哪个拟合更好」，而是\textbf{残差形状}与\textbf{覆盖度区间}：
（1）先看 $V$–$p$ 曲线有无平台——有平台才谈得上 $V_m$；（2）再在双对数坐标（Freundlich）
与 $V$–$\ln p$ 坐标（Temkin）中段看是否直线；（3）比较三种线性化的残差：Langmuir 常在低压
段偏高、Temkin 中段最好、Freundlich 全程都可能「看起来不错」，但其常数没有位点数目含义；
（4）凡要报\textbf{位点数目或分散度}，只能用给出 $V_m$ 的形式（Langmuir 或 BET），并用
5.3 节的脉冲饱和法交叉验证。
\end{dataflow}

\begin{pitfall}
\textbf{在真实催化剂上直接用 Langmuir 是常见错误。} 负载型金属催化剂至少在三个层面违反其
假设：多种晶面与配位不饱和位（吸附热不同）、吸附质之间的横向排斥、强吸附物种阻塞邻近位点；
表现为 $p/V$–$p$ 图弯曲、不同温度下求得的 $V_m$ 不一致。\textbf{另一类错误是把中毒当作
非均匀性}：硫中毒或不可逆强吸附同样让吸附量随压力呈幂律上升，看起来「很 Freundlich」，
但机制与吸附热分布无关——中毒型数据重新抽真空/升温后不可逆，真实非均匀吸附则可逆。
\end{pitfall}

\begin{crosscheck}
等温式给出的 $V_m$ 只说明「有多少可吸附位」，不说明「这些位是否具有催化活性」，必须与
5.3 节的 TPD（位点强度分布）和反应动力学测试（转换频率）互补回查：若 TPD 显示吸附全部
集中在过高温度，则 $V_m$ 表征的是强吸附位而非参与反应的活性位。
此处已对照原页面图像（书内 p.307—310 / PDF p.326—329，§5.2）。原书的符号约定与本稿不同，
逐条记录如下：\textbf{Freundlich} 原书印式 (5-29) $\theta=cp^{1/a}\ (a>1)$、式 (5-32)
$\theta=(a_0p)^{RT/Q_{\text{adm}}}=cp^{1/a}$（其中 $c=a_0^{RT/Q_{\text{adm}}}$、$a=Q_{\text{adm}}/RT$）、
实验表达式 (5-33) $\ln V=\ln V_{\text{m}}+\frac{RT}{Q_{\text{adm}}}\ln a_0+\frac{RT}{Q_{\text{adm}}}\ln p$，
即原书用 \textbf{$c$、$a$}（$a>1$，$a$ 为与吸附物种相互作用有关的常数，$a>1$ 视为吸附分子相互
排斥的结果），本稿用 $K_F$、$n$，且本稿 $V=K_Fp^{1/n}$ 把 $V_{\text{m}}$ 并入常数，属改写。
\textbf{Temkin} 原书印式 (5-34) $Q_{\text{ad}}=Q_{\text{ad}0}(1-\alpha\theta)$、
式 (5-35) $\theta=\frac{RT}{\alpha Q_{\text{ad}0}}\ln(K_0e^{Q_{\text{ad}0}/RT}p)=A_0\ln(B_0p)$、
实验表达式 (5-36) $V=V_{\text{m}}(RT/\alpha Q_{\text{ad}0})(\ln B_0+\ln p)$，即原书用
$Q_{\text{ad}0}$（覆盖度为 0 时的初始吸附热）与 $\alpha$，本稿写 $\Delta H_0-\alpha\theta$（符号与
原书的 $Q_{\text{ad}0}(1-\alpha\theta)$ 等价，$q=-\Delta H$ 故正负号相反），$K_0$ 与原书一致。
\textbf{Langmuir} 原书印 4 条假设（「(1) 吸附只能发生在空吸附位上；(2) 每个吸附位只能吸附一个
分子或原子…达到饱和覆盖度；(3) 吸附热与覆盖度无关，也就是说被吸附分子之间无相互作用；
(4) 吸附和脱附过程一般处于平衡状态」）、式 (5-17)~(5-23)、解离吸附式 (5-24)~(5-27) 及
$n$ 个吸附位作用的式 (5-28) $\theta=(Kp)^{1/n}/[1+(Kp)^{1/n}]$，与本稿一致。另原书 图 5-1 图题
「三种吸附等温式中吸附热 $Q$ 随覆盖度 $\theta$ 的变化：I：Langmuir 等温式；II：Temkin 等温式；
III：Freundlich 等温式」，与本稿一致。
\end{crosscheck}

\subsection{动态分析方法理论}

\subsubsection{程序升温脱附（TPD）基本原理}

\chapref{ch:01} 的容量法与脉冲法都在恒温下测吸附量，TPD 则把\textbf{温度作为扫描变量}。
样品预吸附气体并吹扫掉物理吸附层，再在惰性载气（\ce{He}、\ce{Ar} 或 \ce{N2}）中以恒定
速率升温，下游用热导检测器或质谱监测脱附产物浓度，得到脱附速率对温度的曲线即 TPD 谱。
脱附速率同时依赖温度与覆盖度：升温开始时 $\theta$ 大，速率急剧上升，主要受温度控制；随
吸附分子脱出 $\theta$ 下降，速率转而由 $\theta$ 决定并开始减小；$\theta=0$ 时速率归零。
这一「先指数上升、后衰减」的竞争正是谱峰存在的原因。

\begin{principle}
\textbf{峰面积给出吸附量，峰温度给出吸附强度，峰数目给出位点种类数}——这是 TPD 判读的
全部基础。原书列出的定量信息为：吸附类型（活性位）数目、各类型的强度、每个类型中质点
的数目（活性位密度）、脱附反应级数、表面能量分布；优点是设备简单、几乎覆盖所有实用
催化剂、可原位考察吸附分子与表面的反应
\textsuperscript{\src{S1 §5.3.1，书内 p.311—312 / PDF p.330—331}}。
\end{principle}

\begin{table}[htbp]
\centering
\caption{TPD 谱能给出什么、不能给出什么\src{据 S1 §5.3.1，书内 p.311—312 / PDF p.330—331 归纳。}}
\label{tab:ch05-tpd}
{\srcsize
\begin{tabularx}{\columnwidth}{@{}lXX@{}}
\toprule
信息维度 & 能给出（有前提） & 不能给出（或需外部证据） \\
\midrule
峰数目 & 不同吸附强度物种的\textbf{下限}个数 & 每个峰的化学身份；重叠峰会被读成一个位点 \\
峰面积 & 脱附物种的相对数量；标定后为绝对吸附量 & 未经标定的绝对量；脉冲法标定是前提 \\
峰温度 & 同一系列样品间相对吸附强度的排序 & 直接读出的 $E_d$；须满足 Redhead 前提 \\
峰形/半宽 & 脱附级数线索（一级不对称、二级近对称） & 唯一的级数值；扩散与再吸附同样改变峰形 \\
覆盖度依赖 & 随初始 $\theta$ 的峰温漂移 → 表面非均匀性 & 能量分布的定量形式；需多种 $\theta$ 序列 \\
反应性 & 脱附产物分布（质谱）→ 表面反应路径 & 稳态反应速率；TPD 在无反应物分压条件下进行 \\
\bottomrule
\end{tabularx}}
\end{table}

\emph{核对说明}：已对照原页面图像（书内 p.311—312 / PDF p.330—331）。原书
p.311 列出 TPD 谱可获得的五项定量信息：「①吸附类型（活性中心）的数目；②吸附类型的强度（能量）；
③每个吸附类型中质点的数目（活性中心密度）；④脱附反应级数（吸附质点的相互作用）；⑤表面能量分布（表面均匀性程度）」，并给出 Wigner-Polanyi 式（式 5-37）$N=-V_{\mathrm m}\,\mathrm
d\theta/\mathrm dt=A\theta^{n}\exp[-E_{\mathrm
d}(\theta)/RT]$ 与 TPD 法的主要优点（①设备简单；②研究范围大，几乎可包括所有的实用催化剂；
③原位考察吸附分子和固体表面的反应情况）。p.312 给出热导检测的 TPD 流程（图 5-3 的 12 个部件）
与脉冲进样至峰面积不变的饱和判据（图 5-4）。本表「能给出」一列即上述五项定量信息的转写；
\textbf{「不能给出（或需外部证据）」一列与所列 Redhead 前提（一级脱附、不再吸附、指前因子已知、
升温速率恒定）为本改编本据 §5.3.3 与原书对各条判据的限定归纳，原书未以表格形式逐条列出}。


\subsubsection{实验装置、条件与谱图定性分析}

装置由三个单元组成：气体净化与切换、反应与控温、分析测量。典型流程用两个六通阀：预处理
气体经阀 6 通过样品，达规定温度与时间后切换阀 6 使载气流经样品，再由阀 7 脉冲进样；第一个
脉冲出峰面积最小（甚至不出峰），随脉冲次数增加峰面积增大，直至不变，即认为表面已被吸附质
饱和。以达到饱和时的峰面积为标定峰 $A_0$（此时进气不再被吸附、全部随载气流出），
每次脉冲进气量（标准状态）为 $V_0$，则第 $i$ 次脉冲被吸附的气量为
$(V_0/A_0)(A_0-A_i)$，累加得吸附量
\textsuperscript{\src{S1 §5.3.2，书内 p.313 / PDF p.332}}：
\begin{equation}
V_{\text{吸}} = \frac{V_0}{A_0}\left[\left(A_0 - A_1\right)+\left(A_0 - A_2\right)+\left(A_0 - A_3\right)+\cdots\right]
\label{eq:ch05-pulse}
\end{equation}
式中求和覆盖从第一次脉冲到峰面积达到饱和值 $A_0$ 的那一次脉冲：达到饱和后
$A_i\approx A_0$，后续各项贡献趋于零，故「求和到哪一次」不影响结果，也不会把末次脉冲重复
计入。量纲检查：$[V_0/A_0]\times[A]=\text{体积}$，与 $V_{\text{吸}}$ 同；$A_0$ 与
$A_i$ 用同一检测器增益测得，故 $V_{\text{吸}}$ 与增益无关。
此处已对照原页面图像（书内 p.313 / PDF p.332）：原书式 (5-38) 逐字为
$V=\dfrac{V_0}{A_0}[(A_0-A_1)+(A_0-A_2)+(A_0-A_3)+\cdots]$，是\textbf{开放式展开求和}而非带上下限
的 $\sum$；原书并明确「取达到饱和时的峰面积为标定峰（以 $A_0$ 表示），已知每次脉冲进入的
气体量体积（\textbf{标准状态}）为 $V_0$」——即 $V_0$ 已指明为标准状态体积，本节据此补写
求和上界 $n$ 的写法与量纲检查，属改写而非原书原式。

饱和后继续用载气吹扫至热导基线平衡以脱除物理吸附，再进行程序升温。定性判读三条：
\textbf{峰数目}表征不同吸附强度的物种数目，\textbf{峰面积}表征脱附物种的相对数量，
\textbf{峰温}表征脱附物种在表面的吸附强度。若吸附物与表面发生反应或吸附分子解离，脱附产物
不是单一原始吸附质，此时必须用质谱而非热导检测。

\begin{measure}
为保证 TPD 曲线重复可靠，原书给出的条件区间为：载气流速一般 \qtyrange{30}{50}{mL.min^{-1}}；
装填量一般 \qtyrange{50}{200}{mg}，粒度 \qtyrange{40}{80}{目}；升温速率一般
\qtyrange{10}{20}{K.min^{-1}}，且应有较好的升温线性关系。升温速率过大时峰容易重叠、
造成信息损失；过小时信号减弱且实验时间延长。\textbf{升温线性度是必须记录并报告的量}：
所有基于峰温的动力学计算都假定 $\beta$ 恒定。
\end{measure}

\textbf{纵轴用 $\mathrm{d}\theta/\mathrm{d}T$ 还是 $\mathrm{d}\theta/\mathrm{d}t$，决定峰面积能否
直接标定吸附量。}程序升温实验里检测器记录的是\textbf{按时间}的脱附速率
$\mathrm{d}\theta/\mathrm{d}t$，两者只差一个恒定的升温速率：
$\mathrm{d}\theta/\mathrm{d}T=(\mathrm{d}\theta/\mathrm{d}t)/\beta$。起始覆盖度相同时，
$\mathrm{d}\theta/\mathrm{d}t$ 的峰面积随 $\beta$ 成比例放大；除以 $\beta$ 后的
$\mathrm{d}\theta/\mathrm{d}T$ 峰面积才等于 $\theta_0$、与 $\beta$ 无关
（图~\ref{fig:ch05-tpd}）。所以\textbf{比较不同 $\beta$ 的谱图时只能比峰位、不能比峰高}；
要比峰面积，必须先除 $\beta$，或用同一 $\beta$ 测同一批样品。

\begin{figure}[htbp]
\centering
% 图源：解 Polanyi--Wigner 方程 d(theta)/dt = -A exp(-Ea/RT) theta^n（书内式(5-39) 所用模型），
%       T = T0 + beta t，theta(0) = theta_0 = 1（饱和覆盖），T0 = 300 K。
%       单位（R5-C01）：beta = 5/10/20 K/min；A = 7.345100e+08 s^-1 = 4.407060e+10 min^-1，
%       由 n = 1、beta = 10 K/min、Tp = 570.000 K 的解析峰温条件标定；积分在秒制下完成。
%       纵轴是正脱附速率 -d(theta)/dT = (A/beta_s) exp(-Ea/RT) theta^n（R5-C03 负号约定），
%       故全温区积分 int (-dtheta/dT) dT = theta_0 = 1.000000（与本图图注、正文数字一致）：
%         (a) beta = 5 / 10 / 20 K/min -> 1.000000 / 1.000000 / 1.000000
%         (b) n = 1 / 2（beta = 10 K/min）-> 1.000000 / 1.000000
%       下表坐标由 tools/gen_tpd_table.py 生成（--patch 直接回写本块）。
\begin{tikzpicture}
\begin{groupplot}[
  group style={group size=2 by 1, horizontal sep=7mm},
  width=50mm, height=40mm,
  xmin=400, xmax=900, ymin=0, ymax=0.019,
  xtick={400,500,600,700,800,900}, minor xtick={450,550,650,750,850},
  ytick={0,0.005,0.010,0.015}, yticklabels={0,5,10,15},
  scaled x ticks=false, scaled y ticks=false,   % 只保留 ylabel 里的一套倍率（R5-U01）
  xlabel={$T$ / K}, ylabel={$-\mathrm{d}\theta/\mathrm{d}T$ / $10^{-3}\,\unit{K^{-1}}$},
  tick label style={font=\srcsize}, label style={font=\srcsize},
  title style={font=\srcsize},
  legend style={font=\srcsize, draw=none, fill=none},
  every axis plot/.append style={thick},
]
\nextgroupplot[title={(a) 升温速率（一级）},
  legend style={at={(0.5,-0.50)}, anchor=north, legend columns=1,
                nodes={inner sep=1pt}}]
\addplot[blue, thick] table[row sep=\\, y index=1] {
%%  T/K      -dtheta/dT / K^{-1}
  400.0  1.88401e-06 \\
  410.0  4.54222e-06 \\
  420.0  1.05011e-05 \\
  430.0  2.33470e-05 \\
  440.0  5.00462e-05 \\
  450.0  1.03661e-04 \\
  460.0  2.07858e-04 \\
  470.0  4.04013e-04 \\
  480.0  7.61651e-04 \\
  490.0  1.39195e-03 \\
  500.0  2.46052e-03 \\
  502.0  2.74543e-03 \\
  504.0  3.05852e-03 \\
  506.0  3.40174e-03 \\
  508.0  3.77702e-03 \\
  510.0  4.18623e-03 \\
  512.0  4.63108e-03 \\
  514.0  5.11311e-03 \\
  516.0  5.63357e-03 \\
  518.0  6.19336e-03 \\
  520.0  6.79287e-03 \\
  522.0  7.43194e-03 \\
  524.0  8.10962e-03 \\
  526.0  8.82411e-03 \\
  528.0  9.57257e-03 \\
  530.0  1.03509e-02 \\
  532.0  1.11536e-02 \\
  534.0  1.19738e-02 \\
  536.0  1.28026e-02 \\
  538.0  1.36297e-02 \\
  540.0  1.44424e-02 \\
  542.0  1.52267e-02 \\
  544.0  1.59663e-02 \\
  546.0  1.66435e-02 \\
  548.0  1.72394e-02 \\
  550.0  1.77341e-02 \\
  552.0  1.81075e-02 \\
  554.0  1.83404e-02 \\
  556.0  1.84149e-02 \\
  558.0  1.83158e-02 \\
  560.0  1.80319e-02 \\
  562.0  1.75569e-02 \\
  564.0  1.68909e-02 \\
  566.0  1.60409e-02 \\
  568.0  1.50213e-02 \\
  570.0  1.38546e-02 \\
  572.0  1.25705e-02 \\
  574.0  1.12047e-02 \\
  576.0  9.79745e-03 \\
  578.0  8.39110e-03 \\
  580.0  7.02740e-03 \\
  582.0  5.74464e-03 \\
  584.0  4.57496e-03 \\
  586.0  3.54217e-03 \\
  588.0  2.66040e-03 \\
  590.0  1.93369e-03 \\
  592.0  1.35668e-03 \\
  594.0  9.16292e-04 \\
  596.0  5.93995e-04 \\
  598.0  3.68435e-04 \\
  600.0  2.17927e-04 \\
  602.0  1.22483e-04 \\
  604.0  6.51608e-05 \\
  606.0  3.26787e-05 \\
  608.0  1.53819e-05 \\
  610.0  6.76378e-06 \\
  612.0  2.76466e-06 \\
  614.0  1.04486e-06 \\
  616.0  3.63060e-07 \\
  618.0  1.15287e-07 \\
  620.0  3.32409e-08 \\
  630.0  1.20687e-11 \\
  640.0  1.14507e-16 \\
  650.0  0.00000e+00 \\
  660.0  0.00000e+00 \\
  670.0  0.00000e+00 \\
  680.0  0.00000e+00 \\
  690.0  0.00000e+00 \\
  700.0  0.00000e+00 \\
  710.0  0.00000e+00 \\
  720.0  0.00000e+00 \\
  730.0  0.00000e+00 \\
  740.0  0.00000e+00 \\
  750.0  0.00000e+00 \\
  760.0  0.00000e+00 \\
  770.0  0.00000e+00 \\
  780.0  0.00000e+00 \\
  790.0  0.00000e+00 \\
  800.0  0.00000e+00 \\
  810.0  0.00000e+00 \\
  820.0  0.00000e+00 \\
  830.0  0.00000e+00 \\
  840.0  0.00000e+00 \\
  850.0  0.00000e+00 \\
  860.0  0.00000e+00 \\
  870.0  0.00000e+00 \\
  880.0  0.00000e+00 \\
  890.0  0.00000e+00 \\
  900.0  0.00000e+00 \\
};\addlegendentry{$\beta=5$}
\addplot[red, thick] table[row sep=\\, y index=1] {
%%  T/K      -dtheta/dT / K^{-1}
  400.0  9.42017e-07 \\
  410.0  2.27117e-06 \\
  420.0  5.25087e-06 \\
  430.0  1.16752e-05 \\
  440.0  2.50310e-05 \\
  450.0  5.18663e-05 \\
  460.0  1.04079e-04 \\
  470.0  2.02598e-04 \\
  480.0  3.83032e-04 \\
  490.0  7.03750e-04 \\
  500.0  1.25614e-03 \\
  502.0  1.40541e-03 \\
  504.0  1.57046e-03 \\
  506.0  1.75267e-03 \\
  508.0  1.95348e-03 \\
  510.0  2.17440e-03 \\
  512.0  2.41698e-03 \\
  514.0  2.68282e-03 \\
  516.0  2.97355e-03 \\
  518.0  3.29076e-03 \\
  520.0  3.63605e-03 \\
  522.0  4.01094e-03 \\
  524.0  4.41683e-03 \\
  526.0  4.85498e-03 \\
  528.0  5.32642e-03 \\
  530.0  5.83188e-03 \\
  532.0  6.37174e-03 \\
  534.0  6.94591e-03 \\
  536.0  7.55375e-03 \\
  538.0  8.19391e-03 \\
  540.0  8.86430e-03 \\
  542.0  9.56186e-03 \\
  544.0  1.02825e-02 \\
  546.0  1.10211e-02 \\
  548.0  1.17709e-02 \\
  550.0  1.25242e-02 \\
  552.0  1.32716e-02 \\
  554.0  1.40021e-02 \\
  556.0  1.47035e-02 \\
  558.0  1.53622e-02 \\
  560.0  1.59632e-02 \\
  562.0  1.64908e-02 \\
  564.0  1.69285e-02 \\
  566.0  1.72600e-02 \\
  568.0  1.74694e-02 \\
  570.0  1.75421e-02 \\
  572.0  1.74656e-02 \\
  574.0  1.72306e-02 \\
  576.0  1.68312e-02 \\
  578.0  1.62665e-02 \\
  580.0  1.55410e-02 \\
  582.0  1.46650e-02 \\
  584.0  1.36548e-02 \\
  586.0  1.25326e-02 \\
  588.0  1.13259e-02 \\
  590.0  1.00661e-02 \\
  592.0  8.78720e-03 \\
  594.0  7.52408e-03 \\
  596.0  6.31005e-03 \\
  598.0  5.17498e-03 \\
  600.0  4.14336e-03 \\
  602.0  3.23287e-03 \\
  604.0  2.45348e-03 \\
  606.0  1.80737e-03 \\
  608.0  1.28953e-03 \\
  610.0  8.89040e-04 \\
  612.0  5.90796e-04 \\
  614.0  3.77421e-04 \\
  616.0  2.31131e-04 \\
  618.0  1.35277e-04 \\
  620.0  7.54274e-05 \\
  630.0  1.72886e-06 \\
  640.0  6.37182e-09 \\
  650.0  1.87666e-12 \\
  660.0  1.74736e-17 \\
  670.0  0.00000e+00 \\
  680.0  0.00000e+00 \\
  690.0  0.00000e+00 \\
  700.0  0.00000e+00 \\
  710.0  0.00000e+00 \\
  720.0  0.00000e+00 \\
  730.0  0.00000e+00 \\
  740.0  0.00000e+00 \\
  750.0  0.00000e+00 \\
  760.0  0.00000e+00 \\
  770.0  0.00000e+00 \\
  780.0  0.00000e+00 \\
  790.0  0.00000e+00 \\
  800.0  0.00000e+00 \\
  810.0  0.00000e+00 \\
  820.0  0.00000e+00 \\
  830.0  0.00000e+00 \\
  840.0  0.00000e+00 \\
  850.0  0.00000e+00 \\
  860.0  0.00000e+00 \\
  870.0  0.00000e+00 \\
  880.0  0.00000e+00 \\
  890.0  0.00000e+00 \\
  900.0  0.00000e+00 \\
};\addlegendentry{$\beta=10$}
\addplot[black, thick] table[row sep=\\, y index=1] {
%%  T/K      -dtheta/dT / K^{-1}
  400.0  4.71011e-07 \\
  410.0  1.13560e-06 \\
  420.0  2.62551e-06 \\
  430.0  5.83800e-06 \\
  440.0  1.25175e-05 \\
  450.0  2.59421e-05 \\
  460.0  5.20769e-05 \\
  470.0  1.01447e-04 \\
  480.0  1.92070e-04 \\
  490.0  3.53835e-04 \\
  500.0  6.34643e-04 \\
  502.0  7.11023e-04 \\
  504.0  7.95737e-04 \\
  506.0  8.89579e-04 \\
  508.0  9.93399e-04 \\
  510.0  1.10811e-03 \\
  512.0  1.23468e-03 \\
  514.0  1.37414e-03 \\
  516.0  1.52758e-03 \\
  518.0  1.69616e-03 \\
  520.0  1.88106e-03 \\
  522.0  2.08355e-03 \\
  524.0  2.30489e-03 \\
  526.0  2.54642e-03 \\
  528.0  2.80946e-03 \\
  530.0  3.09534e-03 \\
  532.0  3.40537e-03 \\
  534.0  3.74079e-03 \\
  536.0  4.10278e-03 \\
  538.0  4.49241e-03 \\
  540.0  4.91056e-03 \\
  542.0  5.35792e-03 \\
  544.0  5.83490e-03 \\
  546.0  6.34159e-03 \\
  548.0  6.87766e-03 \\
  550.0  7.44229e-03 \\
  552.0  8.03411e-03 \\
  554.0  8.65108e-03 \\
  556.0  9.29039e-03 \\
  558.0  9.94840e-03 \\
  560.0  1.06205e-02 \\
  562.0  1.13011e-02 \\
  564.0  1.19836e-02 \\
  566.0  1.26599e-02 \\
  568.0  1.33213e-02 \\
  570.0  1.39576e-02 \\
  572.0  1.45575e-02 \\
  574.0  1.51089e-02 \\
  576.0  1.55991e-02 \\
  578.0  1.60146e-02 \\
  580.0  1.63420e-02 \\
  582.0  1.65682e-02 \\
  584.0  1.66809e-02 \\
  586.0  1.66691e-02 \\
  588.0  1.65242e-02 \\
  590.0  1.62398e-02 \\
  592.0  1.58133e-02 \\
  594.0  1.52457e-02 \\
  596.0  1.45426e-02 \\
  598.0  1.37141e-02 \\
  600.0  1.27749e-02 \\
  602.0  1.17444e-02 \\
  604.0  1.06455e-02 \\
  606.0  9.50434e-03 \\
  608.0  8.34884e-03 \\
  610.0  7.20729e-03 \\
  612.0  6.10689e-03 \\
  614.0  5.07218e-03 \\
  616.0  4.12366e-03 \\
  618.0  3.27666e-03 \\
  620.0  2.54064e-03 \\
  630.0  4.62693e-04 \\
  640.0  3.35950e-05 \\
  650.0  6.85760e-07 \\
  660.0  2.46950e-09 \\
  670.0  8.47768e-13 \\
  680.0  1.24784e-17 \\
  690.0  0.00000e+00 \\
  700.0  0.00000e+00 \\
  710.0  0.00000e+00 \\
  720.0  0.00000e+00 \\
  730.0  0.00000e+00 \\
  740.0  0.00000e+00 \\
  750.0  0.00000e+00 \\
  760.0  0.00000e+00 \\
  770.0  0.00000e+00 \\
  780.0  0.00000e+00 \\
  790.0  0.00000e+00 \\
  800.0  0.00000e+00 \\
  810.0  0.00000e+00 \\
  820.0  0.00000e+00 \\
  830.0  0.00000e+00 \\
  840.0  0.00000e+00 \\
  850.0  0.00000e+00 \\
  860.0  0.00000e+00 \\
  870.0  0.00000e+00 \\
  880.0  0.00000e+00 \\
  890.0  0.00000e+00 \\
  900.0  0.00000e+00 \\
};\addlegendentry{$\beta=20$}
\nextgroupplot[title={(b) 脱附级数（$\beta=10$）},
  legend style={at={(0.5,-0.50)}, anchor=north, legend columns=1,
                nodes={inner sep=1pt}}]
\addplot[red, thick] table[row sep=\\, y index=1] {
%%  T/K      -dtheta/dT / K^{-1}
  400.0  9.42017e-07 \\
  410.0  2.27117e-06 \\
  420.0  5.25087e-06 \\
  430.0  1.16752e-05 \\
  440.0  2.50310e-05 \\
  450.0  5.18663e-05 \\
  460.0  1.04079e-04 \\
  470.0  2.02598e-04 \\
  480.0  3.83032e-04 \\
  490.0  7.03750e-04 \\
  500.0  1.25614e-03 \\
  502.0  1.40541e-03 \\
  504.0  1.57046e-03 \\
  506.0  1.75267e-03 \\
  508.0  1.95348e-03 \\
  510.0  2.17440e-03 \\
  512.0  2.41698e-03 \\
  514.0  2.68282e-03 \\
  516.0  2.97355e-03 \\
  518.0  3.29076e-03 \\
  520.0  3.63605e-03 \\
  522.0  4.01094e-03 \\
  524.0  4.41683e-03 \\
  526.0  4.85498e-03 \\
  528.0  5.32642e-03 \\
  530.0  5.83188e-03 \\
  532.0  6.37174e-03 \\
  534.0  6.94591e-03 \\
  536.0  7.55375e-03 \\
  538.0  8.19391e-03 \\
  540.0  8.86430e-03 \\
  542.0  9.56186e-03 \\
  544.0  1.02825e-02 \\
  546.0  1.10211e-02 \\
  548.0  1.17709e-02 \\
  550.0  1.25242e-02 \\
  552.0  1.32716e-02 \\
  554.0  1.40021e-02 \\
  556.0  1.47035e-02 \\
  558.0  1.53622e-02 \\
  560.0  1.59632e-02 \\
  562.0  1.64908e-02 \\
  564.0  1.69285e-02 \\
  566.0  1.72600e-02 \\
  568.0  1.74694e-02 \\
  570.0  1.75421e-02 \\
  572.0  1.74656e-02 \\
  574.0  1.72306e-02 \\
  576.0  1.68312e-02 \\
  578.0  1.62665e-02 \\
  580.0  1.55410e-02 \\
  582.0  1.46650e-02 \\
  584.0  1.36548e-02 \\
  586.0  1.25326e-02 \\
  588.0  1.13259e-02 \\
  590.0  1.00661e-02 \\
  592.0  8.78720e-03 \\
  594.0  7.52408e-03 \\
  596.0  6.31005e-03 \\
  598.0  5.17498e-03 \\
  600.0  4.14336e-03 \\
  602.0  3.23287e-03 \\
  604.0  2.45348e-03 \\
  606.0  1.80737e-03 \\
  608.0  1.28953e-03 \\
  610.0  8.89040e-04 \\
  612.0  5.90796e-04 \\
  614.0  3.77421e-04 \\
  616.0  2.31131e-04 \\
  618.0  1.35277e-04 \\
  620.0  7.54274e-05 \\
  630.0  1.72886e-06 \\
  640.0  6.37182e-09 \\
  650.0  1.87666e-12 \\
  660.0  1.74736e-17 \\
  670.0  0.00000e+00 \\
  680.0  0.00000e+00 \\
  690.0  0.00000e+00 \\
  700.0  0.00000e+00 \\
  710.0  0.00000e+00 \\
  720.0  0.00000e+00 \\
  730.0  0.00000e+00 \\
  740.0  0.00000e+00 \\
  750.0  0.00000e+00 \\
  760.0  0.00000e+00 \\
  770.0  0.00000e+00 \\
  780.0  0.00000e+00 \\
  790.0  0.00000e+00 \\
  800.0  0.00000e+00 \\
  810.0  0.00000e+00 \\
  820.0  0.00000e+00 \\
  830.0  0.00000e+00 \\
  840.0  0.00000e+00 \\
  850.0  0.00000e+00 \\
  860.0  0.00000e+00 \\
  870.0  0.00000e+00 \\
  880.0  0.00000e+00 \\
  890.0  0.00000e+00 \\
  900.0  0.00000e+00 \\
};\addlegendentry{$n=1$，拖尾}
\addplot[black, dashed, thick] table[row sep=\\, y index=1] {
%%  T/K      -dtheta/dT / K^{-1}
  400.0  9.42007e-07 \\
  410.0  2.27111e-06 \\
  420.0  5.25055e-06 \\
  430.0  1.16735e-05 \\
  440.0  2.50231e-05 \\
  450.0  5.18307e-05 \\
  460.0  1.03929e-04 \\
  470.0  2.02008e-04 \\
  480.0  3.80838e-04 \\
  490.0  6.96062e-04 \\
  500.0  1.23079e-03 \\
  502.0  1.37347e-03 \\
  504.0  1.53032e-03 \\
  506.0  1.70237e-03 \\
  508.0  1.89063e-03 \\
  510.0  2.09608e-03 \\
  512.0  2.31968e-03 \\
  514.0  2.56232e-03 \\
  516.0  2.82477e-03 \\
  518.0  3.10770e-03 \\
  520.0  3.41157e-03 \\
  522.0  3.73666e-03 \\
  524.0  4.08297e-03 \\
  526.0  4.45020e-03 \\
  528.0  4.83766e-03 \\
  530.0  5.24431e-03 \\
  532.0  5.66859e-03 \\
  534.0  6.10850e-03 \\
  536.0  6.56148e-03 \\
  538.0  7.02443e-03 \\
  540.0  7.49374e-03 \\
  542.0  7.96522e-03 \\
  544.0  8.43422e-03 \\
  546.0  8.89566e-03 \\
  548.0  9.34411e-03 \\
  550.0  9.77393e-03 \\
  552.0  1.01794e-02 \\
  554.0  1.05547e-02 \\
  556.0  1.08945e-02 \\
  558.0  1.11937e-02 \\
  560.0  1.14476e-02 \\
  562.0  1.16526e-02 \\
  564.0  1.18055e-02 \\
  566.0  1.19043e-02 \\
  568.0  1.19480e-02 \\
  570.0  1.19364e-02 \\
  572.0  1.18706e-02 \\
  574.0  1.17523e-02 \\
  576.0  1.15845e-02 \\
  578.0  1.13707e-02 \\
  580.0  1.11149e-02 \\
  582.0  1.08220e-02 \\
  584.0  1.04967e-02 \\
  586.0  1.01442e-02 \\
  588.0  9.76968e-03 \\
  590.0  9.37818e-03 \\
  592.0  8.97452e-03 \\
  594.0  8.56324e-03 \\
  596.0  8.14853e-03 \\
  598.0  7.73416e-03 \\
  600.0  7.32347e-03 \\
  602.0  6.91935e-03 \\
  604.0  6.52427e-03 \\
  606.0  6.14026e-03 \\
  608.0  5.76895e-03 \\
  610.0  5.41161e-03 \\
  612.0  5.06916e-03 \\
  614.0  4.74225e-03 \\
  616.0  4.43124e-03 \\
  618.0  4.13629e-03 \\
  620.0  3.85733e-03 \\
  630.0  2.69163e-03 \\
  640.0  1.85816e-03 \\
  650.0  1.27815e-03 \\
  660.0  8.80104e-04 \\
  670.0  6.08440e-04 \\
  680.0  4.23081e-04 \\
  690.0  2.96222e-04 \\
  700.0  2.08956e-04 \\
  710.0  1.48545e-04 \\
  720.0  1.06430e-04 \\
  730.0  7.68510e-05 \\
  740.0  5.59197e-05 \\
  750.0  4.09954e-05 \\
  760.0  3.02743e-05 \\
  770.0  2.25158e-05 \\
  780.0  1.68609e-05 \\
  790.0  1.27104e-05 \\
  800.0  9.64323e-06 \\
  810.0  7.36178e-06 \\
  820.0  5.65389e-06 \\
  830.0  4.36747e-06 \\
  840.0  3.39270e-06 \\
  850.0  2.64980e-06 \\
  860.0  2.08044e-06 \\
  870.0  1.64171e-06 \\
  880.0  1.30187e-06 \\
  890.0  1.03728e-06 \\
  900.0  8.30270e-07 \\
};\addlegendentry{$n=2$，对称}
\end{groupplot}
\end{tikzpicture}
\caption{TPD 谱的模型计算：\textbf{峰面积标定吸附量，峰位标定动力学}。曲线由
Polanyi--Wigner 方程数值求解，取相同初始覆盖度 $\theta_0=1$（饱和）、
$E_d=\qty{120}{kJ.mol^{-1}}$、$A=\qty{7.35e8}{s^{-1}}$（$=\qty{4.41e10}{min^{-1}}$；
由 $n=1$、$\beta=\qty{10}{K.min^{-1}}$、$T_p=\qty{570.000}{K}$ 的解析峰温条件标定）、
$T_0=\qty{300}{K}$。纵轴为正脱附速率 $-\mathrm{d}\theta/\mathrm{d}T$，即把按时间计的
脱附速率 $-\mathrm{d}\theta/\mathrm{d}t$ 除以 $\beta$：脱附使 $\theta$ 随 $T$ 下降，加负号后
曲线为正，与式~\eqref{eq:ch05-peak} 的导数符号约定一致。\textbf{(a)} 一级脱附，$\beta$ 由 $5$ 增到
$\qty{20}{K.min^{-1}}$：峰温由 $556$ 右移到 $\qty{585}{K}$（约 $\qty{29}{K}$），
而三条曲线下的面积\textbf{完全相同}
（$\int(-\mathrm{d}\theta/\mathrm{d}T)\,\mathrm{d}T=\theta_0=1.0000$）——
\textbf{面积并不随 $\beta$ 成倍放大}，这正是「峰面积标定吸附量」的依据；
按时间计的速率 $-\mathrm{d}\theta/\mathrm{d}t$ 才随 $\beta$ 成比例放大，两者只差因子
$\beta$，故比较升温速率只能看峰位、不能看峰高。\textbf{(b)} 固定
$\beta=\qty{10}{K.min^{-1}}$ 只改脱附级数：$n=2$ 的峰比 $n=1$ \textbf{更低更宽}——
峰高比 $0.681$，半高全宽由 $52.3$ 增至 $\qty{73.5}{K}$（约 $1.4$ 倍）——
峰位略向低温移（$568.6$ 对 $\qty{570.0}{K}$），峰面积仍是 $1.0000$
（初始覆盖度相同）。局限：只含一级与二级脱附、单一均匀表面与恒定 $\beta$，
未含再吸附、扩散与检测器响应，不可用于直接读数值。横轴两端之外仍有极小尾部
（$\theta$ 的残余 $<10^{-6}$），上列积分按 $300$--$\qty{1500}{K}$ 全温区算得。\modelmark
\figsrc{据 S1 §5.3.1—5.3.3 原理自绘，书内 p.311—315 / PDF p.330—334；坐标表由
\texttt{tools/gen\_tpd\_table.py} 生成}}
\label{fig:ch05-tpd}
\end{figure}

\begin{pitfall}
\textbf{峰重叠。} 相邻位点的峰在半宽可比于峰间距时合并为一个宽峰，被读成「单一均匀位点」。
识别办法是改变升温速率或初始覆盖度：真实两峰的条件响应不同，单峰只整体平移。

\textbf{再吸附抬高峰温。} 载气流速过低、床层过厚或样品量过大时，脱附分子在床层内被重新
吸附，表现为峰温偏高、峰形拖尾，由峰温算得的 $E_d$ 系统性偏大。判据是把载气流速提高一档
重测：真实峰位基本不动，再吸附控制的峰位明显左移。

\textbf{载气流速与升温速率耦合。} 两者共同决定脱附物在检测器处的驻留时间与峰宽，单独优化
其一会得到假的最优条件；报告时必须同时给出 $\beta$、流速、装量和粒度。

\textbf{定量必须标定。} 热导或质谱信号是浓度而非摩尔量，峰面积只有经已知体积脉冲标定
（式~\eqref{eq:ch05-pulse}）后才能换算为吸附量；未标定就比较不同批次样品的峰面积，在
检测器灵敏度漂移时会得出错误的位点数目结论。
\end{pitfall}

\subsubsection{TPD 动力学参数的确定}

Cvetanović 和 Amenorinya 为简化理论处理提出四条假设：\textbf{不发生再吸附}、\textbf{脱附为
一级过程}、\textbf{表面均匀}、\textbf{不存在扩散}。此时脱附速率
$N = -\mathrm{d}\theta/\mathrm{d}t = k_d\,\theta^{\,n}$，$k_d$ 取式~\eqref{eq:ch05-kd}；
设线性升温 $T=T_0+\beta t$，峰顶处 $\mathrm{d}^2\theta/\mathrm{d}T^2=0$，解得峰温条件
\begin{equation}
\frac{E_d}{R\,T_m^{2}} = \frac{A}{\beta}\; n\,\theta_m^{\,n-1}
\exp\!\left(-\frac{E_d}{R\,T_m}\right)
\label{eq:ch05-peak}
\end{equation}
式中下标 $m$ 表示峰最高点处的量，其来源是 $\mathrm{d}\theta/\mathrm{d}T =
-(A/\beta)\,\theta^{\,n}\exp(-E_d/RT)$ 在峰顶处满足 $\mathrm{d}^2\theta/\mathrm{d}T^2=0$。
对一级脱附（$n=1$）$\theta_m$ 项消失，即得 Redhead 方程
\begin{equation}
\frac{E_d}{R\,T_m^{2}} = \frac{A}{\beta}\exp\!\left(-\frac{E_d}{R\,T_m}\right)
\label{eq:ch05-redhead}
\end{equation}
此处已对照原页面图像（书内 p.315 / PDF p.334）。原书不按 $n$ 阶通式起手，而是自
一级脱附出发，在峰顶处令 $\mathrm{d}C/\mathrm{d}T=0$，即
$\dfrac{\partial\theta}{\partial T}=\dfrac{V_sV_mA}{F}\cdot\dfrac{\mathrm{d}(\theta\mathrm{e}^{-E_d/RT})}{\mathrm{d}T}=0$，
从而 $\mathrm{d}(\theta\mathrm{e}^{-E_d/RT})/\mathrm{d}T=0$（式 5-52）；解此微分方程得
$(k_d)_M=\beta E_d/RT_M^2$（式 5-53），再由 $k_d=A\exp(-E_d/RT)$ 得
$(k_d)_M=A\exp(-E_d/RT_M)=\beta E_d/RT_M^2$（式 5-54）。原书并注明「这里 $M$ 指 TPD 峰的
最高点」。因此本稿的 $n\theta_m^{\,n-1}$ 因子与「先写 $n$ 阶、再退化到 $n=1$」的叙述层次
属本改编本为便于讨论非一级脱附而加，原书通篇按一级处理；两者在 $n=1$ 时一致。
原书同处另写「在已知加热速率 $\beta$ 时，假定指前因子 $A$ 也已知，从式 (5-54) 可看出，
脱附活化能和 $T_M$ 具有对应关系」，即该式在原书中是逐温度对应的隐式关系，本稿写成
$E_d/RT_m^{2}=(A/\beta)\exp(-E_d/RT_m)$ 属等价变形。原书未出现常数 3.64 的形式。

式~\eqref{eq:ch05-redhead} 是把\textbf{峰温}翻译成\textbf{脱附活化能}的唯一桥梁，它无解析
解、实用中按假定的 $A$ 值迭代求解；其可换算性完全取决于表~\ref{tab:ch05-redhead} 的四个前提。

\begin{table}[htbp]
\centering
\caption{由峰温换算脱附活化能的四个前提\src{前两条为 S1 §5.3.3 原书明确列出的简化假设（书内 p.313 / PDF p.332），后两条据书内 p.313—316 归纳。}}
\label{tab:ch05-redhead}
{\srcsize
\begin{tabularx}{\columnwidth}{@{}lXX@{}}
\toprule
前提 & 要求 & 被违反时的表现 \\
\midrule
一级脱附 & $n=1$，$\theta_m$ 项消失 & 二级脱附下 $E_d$ 随初始覆盖度变化 \\
不再吸附 & 原书简化假设①，正文式 (5-39)—(5-41) 已略去再吸附项 $k_aC(1-\theta)$ & 峰温偏高、拖尾，$E_d$ 系统性偏大 \\
指前因子已知 & 需假定 $A$（常见 $10^{13}\,\si{s^{-1}}$ 量级） & $E_d$ 与 $\ln A$ 强耦合：$A$ 差 2 个量级则 $E_d$ 差 $4.606\,RT_m$（$T_m=500$ K 时约 $19\,\si{kJ mol^{-1}}$）；本书表 5-3 自身算例的 $A$ 仅 $10^{4}\sim10^{5}\,\si{s^{-1}}$ \\
升温速率已知且恒定 & $\beta$ 线性、并在报告中给出 & 峰位漂移，变升温速率法斜率失效 \\
\bottomrule
\end{tabularx}}
\end{table}

\emph{核对说明}：已对照原页面图像（书内 p.313、p.315—316 / PDF p.332、p.334—335）。本表前两条（「一级脱附」「不再吸附」）对应原书 §5.3.3 开篇逐字列出的 Cvetanović 和 Amenomiya（1967）
四条简化假设中的假设①②，后两条（「指前因子已知」「升温速率已知且恒定」）系据书内 p.313—316 归纳，原书未单独列成表格（详见下文）。表中「$A$ 差 2 个量级则 $E_d$ 差 $4.606\,RT_m$」一句为改编本按 Redhead 近似式 $E_d=RT_m[\ln(AT_m/\beta)-3.64]$
对 $\ln A$ 求差所得：$\ln 100=4.606$，故 $\Delta E_d=R T_m\ln 100=4.606\,R T_m$（$T_m=500$ K 时约 $19\,\si{kJ mol^{-1}}$），原书未作此讨论。原稿此处写作「$E_d$ 约差 10\%」，该比例随 $E_d$ 水平漂移（$E_d=200$ 与 $100\,\si{kJ mol^{-1}}$ 时分别约 10\% 与 19\%），故改为给绝对量。
原书表 5-3 自身算例所印 $A$ 仅 $10^{4}\sim10^{5}\,\si{s^{-1}}$，与该前提并不自洽，
原书亦未就此展开讨论，本稿如实记录、不代为改动（逐值见下文）。

原书 §5.3.3 开篇（书内 p.313 / PDF p.332）逐字列出 Cvetanović 和 Amenomiya 于 1967 年提出的
四条简化假设：「①脱附过程不发生再吸附；②脱附为一级脱附过程；③固体表面是均匀的；
④不存在扩散过程」（假设①对应上表「不再吸附」，假设②对应「一级脱附」，假设③④为表面均匀
与无扩散，本稿未单列）。原书自式的起点为式 (5-39) $N=-\mathrm{d}\theta/\mathrm{d}t=k_d\theta-k_aC(1-\theta)$、
式 (5-40) $-V_m\mathrm{d}\theta/\mathrm{d}t=V_mk_d\theta-k_aC(1-\theta)$、
式 (5-41) $FC=V_s[V_mk_d\theta-k_aC(1-\theta)]$；本稿正文的
$N=-\mathrm{d}\theta/\mathrm{d}t=k_d\theta^{\,n}$ 是把式 (5-39) 在假设①下略去再吸附项、
并按本稿讨论非一级情形的需要改写为 $n$ 阶形式，原书只取一级。

\begin{measure}
\textbf{变升温速率法（不假定 $A$）。} 若不作 $A$ 的假设，可由式~\eqref{eq:ch05-redhead}
整理出：以 $2\ln T_m - \ln\beta$ 对 $1/T_m$ 作图，斜率为 $E_d/R$，由此得到脱附活化能。
实验上需在多个 $\beta$ 下各测一条 TPD 曲线、读出各自的 $T_m$ 后作图。原书给出烯烃在
氧化铝中心上脱附的实例（\ce{C2H4}、\ce{C3H6}、\ce{C4H8}），所得 $E_d$ 与 $A$ 列于其表 5-3；
作者同时指出该法实验工作量较大，并提到用谱线重叠法从单条 TPD 谱线求 $E_d$ 与 $A$ 的
替代途径\textsuperscript{\src{S1 §5.3.3，书内 p.315—316 / PDF p.334—335}}。
此处已对照原页面图像（书内 p.315—316 / PDF p.334—335）。原书式
(5-53) $(k_d)_M=\beta E_d/RT_M^2$、式 (5-54) $(k_d)_M=A\exp(-E_d/RT_M)=\beta E_d/RT_M^2$，
下标 $M$ 指 TPD 峰最高点；式 (5-55) 写作
$2\ln T_M-\ln\beta=E_d/RT_M+\ln(E_d/AR)$，即截距并非零而是 $\ln(E_d/AR)$，
本稿以「斜率即 $E_d/R$」概括、略去截距项，属改写。原书表 5-3「烯烃在 \ce{Al2O3} 上
脱附活化能（kcal/mol）和指前因子（$A$）」逐值核对：\ce{C2H4} $E_d=26.8$、$A=1.6\times10^5$；
\ce{C3H6} $E_d=14.2$、$A=6.8\times10^5$；\ce{C4H8} $E_d=12.1$、$A=1.2\times10^4$
（$E_d$ 单位确为 kcal/mol）。三组 $A$ 值确如原书所印、并非 OCR 失真，且远低于
常用假定 $10^{13}\,\si{s^{-1}}$——即\textbf{原书自身的算例并不满足其表 5-3 前提中的
「指前因子已知且合理」一条}，此处如实记录原书数值，不代为改动。
另据 p.316（PDF p.335）原文，该法「需测定不同升温速率 $\beta$ 时的峰顶温度 $T_M$，
用作图法求脱附活化能 $E_d$ 和指前因子 $A$，因此实验工作量较大」，并提到杨上闰用
谱线重叠法由一条 TPD 谱线求得 $E_d$ 与 $A$（文献[6]）。
\end{measure}

\subsubsection{程序升温还原（TPR）原理}

等温条件下的还原过程一般用成核模型和球收缩模型解释。以球形金属氧化物与 \ce{H2} 反应生成
金属和 \ce{H2O} 为例：\textbf{成核模型}认为反应从成核开始，随核长大与新核形成，金属核与
氧化物之间的反应界面增加、速率加快；核相互接触后界面变小、速率减慢，故 $\alpha$–$t$ 曲线
呈 S 形、$\mathrm{d}\alpha/\mathrm{d}t$ 先增后减。\textbf{球收缩模型}认为迅速成核形成很薄的
金属层后界面不断缩小、速率单调下降；两者区别在于成核速度。

程序升温条件下还原写作 $\ce{G + S -> P}$（\ce{G} 为 \ce{H2}，\ce{S} 为氧化物，\ce{P} 为
产物），速率与 Arrhenius 温度依赖、线性升温的关系为
\begin{equation}
r = \frac{\mathrm{d}\alpha}{\mathrm{d}t} = k\,f(\alpha),\qquad
k = A\exp\!\left(-\frac{E_a}{RT}\right),\qquad
\beta = \frac{\mathrm{d}T}{\mathrm{d}t}
\label{eq:ch05-tpr}
\end{equation}
式中 $\alpha$ 为还原分数。因 $\mathrm{d}/\mathrm{d}t = \beta\,\mathrm{d}/\mathrm{d}T$，
反应与气体浓度无关时有
$\mathrm{d}\alpha/\mathrm{d}T = (A/\beta)(1-\alpha)^{q}\exp(-E_a/RT)$。峰顶处
$\mathrm{d}[G]/\mathrm{d}T = 0$，整理、取对数并假定固体反应级数 $q=1$，得
\begin{equation}
2\ln T_m - \ln\beta + p\ln[G]_M + (q-1)\ln[S]_M = \frac{E_a}{R\,T_m} + \text{常数}
\label{eq:ch05-tpr-peak}
\end{equation}
式中 $p$、$q$ 分别为气体与固体的反应级数，$[G]_M$ 为峰顶处还原气浓度，$[S]_M$ 为峰顶处
固体浓度。已知 $p$、$q$ 时以式~\eqref{eq:ch05-tpr-peak} 左边对 $1/T_m$ 作图，斜率即 $E_a/R$。
由于固体反应级数通常为分数且该式对 $q$ 不敏感，原书进而\textbf{假定 $q=1$}，得化简式
\begin{equation}
2\ln T_m - \ln\beta + p\ln[G]_M = \frac{E_a}{R\,T_m} + \text{常数}
\label{eq:ch05-tpr-peak-q1}
\end{equation}
此处已对照原页面图像（书内 p.319 / PDF p.338）：原书式 (5-70)
$0=E/RT_M^2+q\,\mathrm{d}[S]_M/([S]_M\mathrm{d}T)$、式 (5-71)
$E/RT_M^2=[A[G]_M^p q[S]_M^{q-1}\exp(-E/RT)]/\beta$、式 (5-72) 即上式、
式 (5-73) 即化简式，逐字一致。原书另注：对 \ce{CuO} 还原发现 $p=0$ 时 $E$ 与
\ce{H2} 浓度有关，而 $p=1$ 时 $E$ 与浓度无关。

原书指出该法有两个实务限制。其一，\textbf{它并不实用}：实验拿到的 TPR 信息大多是还原气
（如 \ce{H2}）浓度与温度的关系、谱图由几个峰组成，而非 $\alpha$–$T$ 曲线本身；对流动体系
还原过程表现为 $\Delta[G] = [G]_{\text{进}} - [G]_{\text{出}}$，假定活塞流时低转化率下可由
$r = f\,[G]x$ 描述。其二，\textbf{氢溢流}：在某一物质表面形成的活性物种转移到在相同条件下
自身不能吸附或不能形成活性物种的另一物质表面。多组分催化剂中易还原的氧化物先在较低温度
还原成金属，若该金属能解离活化 \ce{H2}，原子氢便经金属—氧化物界面溢流到氧化物表面使之
还原；原子氢还原性强于分子氢，故存在氢溢流时氧化物还原温度明显降低。\ce{Pt}、\ce{Pd}、
\ce{Ag}、\ce{Cu}、\ce{Ni} 等均有此功能，贵金属作用尤其突出。

\begin{crosscheck}
氢溢流使 TPR 图谱难以按单一还原步骤归属：一个低温峰可能不代表该物种本身易还原，而代表
邻近金属提供了溢流通道。因此\textbf{TPR 的峰归属必须与 XRD、\ce{H2} 消耗量（定量）以及
\ce{CO}-TPR 回查}。原书明确指出，为消除氢溢流影响可选用溢流作用较小的 \ce{CO} 替代
\ce{H2} 作还原气，即 \ce{CO}-TPR\textsuperscript{\src{S1 §5.3.4，书内 p.316—319 /
PDF p.335—338}}。
\end{crosscheck}

\subsubsection{程序升温氧化（TPO）原理}

程序升温氧化的原理与 TPR 基本类似，是在一定升温速率下用氧化性气体（如 \ce{O2}）对催化剂
及其表面物种进行氧化的过程，主要用于研究催化剂表面积炭物种的生成机理。

\begin{principle}
与 TPD 的对应关系可以直接迁移：TPO 谱的峰面积对应\textbf{可氧化物种（积炭）的总量}，
峰温对应\textbf{该物种被氧化的难易程度}（积炭越石墨化、越难氧化，峰温越高），
峰数目对应\textbf{积炭物种的类别数}。因此 TPO 与 TPD 的全部定量前提——标定、载气流速、
升温速率、峰重叠——同样适用，且因为反应产物是 \ce{CO2}/\ce{CO}，需用质谱同时监测
两个通道以避免把部分氧化产物计漏。
\end{principle}

\srcnote{本节（5.1—5.3）据 S1 上册第 5 章，书内 p.301—319 / PDF p.320—338。化学吸附应用、
活性位与分散度测定、TPD/TPR/TPO 在具体体系中的应用见本章后续分片。原文中的
图 5-2—5-7 为原书谱图与装置流程图，因无图像核对能力未裁图，本分片仅以
\modelmark 模型计算图替代。}
